\section{Un échéancier certain}
\label{sec:U2}

L'unité précédente traitait un flux isolé. On ajoute ici le premier ingrédient :
\emph{plusieurs dates}. Rien de conceptuellement neuf --- on additionne des
valeurs actuelles --- mais cette addition est la forme sous laquelle se
présentent la quasi-totalité des engagements d'assurance.

\subsection{La notion : valeur actuelle nette d'un échéancier}

\begin{notion}{Sommer des flux actualisés}
Un \emph{échéancier}\index{échéancier} est une suite de flux $(C_1,\dots,C_n)$
associés à des dates $(t_1,\dots,t_n)$. Sa valeur actuelle est la somme des
valeurs actuelles de chaque flux :
\[
  V_0 \;=\; \sum_{k=1}^{n} \frac{C_k}{(1+r)^{t_k}}.
\]
Lorsque l'échéancier comprend un décaissement initial $C_0$, on parle de
\emph{valeur actuelle nette}\index{valeur actuelle nette} :
\[
  \mathrm{VAN} \;=\; -\,C_0 \;+\; \sum_{k=1}^{n} \frac{C_k}{(1+r)^{t_k}}.
\]
Une opération est créatrice de valeur si sa VAN est positive.
\end{notion}

Deux conventions doivent être fixées avant tout calcul, et elles le seront de
même dans les projections prudentielles : les flux sont-ils supposés survenir en
début ou en fin de période, et le taux est-il constant sur tout l'horizon ? Sauf
mention contraire, on retient ici des flux de fin de période et un taux constant.

\begin{center}
\begin{tikzpicture}[x=1.25cm, y=0.85cm]
  \draw[timeaxis] (-0.4,0) -- (5.2,0);
  \foreach \x/\l in {0/0, 1/1, 2/2, 3/3, 4/4}
    { \node[timenode] at (\x,0) {}; \node[etiq, below=3pt] at (\x,0) {\l}; }
  \draw[fluxneg] (0,0) -- (0,-1.1);
  \node[etiq, below] at (0,-1.15) {$C_0$};
  \draw[flux] (1,0) -- (1,0.75); \node[etiq, above] at (1,0.75) {$C_1$};
  \draw[flux] (2,0) -- (2,1.05); \node[etiq, above] at (2,1.05) {$C_2$};
  \draw[flux] (3,0) -- (3,1.35); \node[etiq, above] at (3,1.35) {$C_3$};
  \draw[flux] (4,0) -- (4,1.6);  \node[etiq, above] at (4,1.6)  {$C_4$};
\end{tikzpicture}
\end{center}

\begin{exercice}[valeur actuelle d'un échéancier]
Un projet génère les flux suivants, en fin d'année : $12\,000$~\euro,
$15\,000$~\euro, $18\,000$~\euro{} et $20\,000$~\euro. Le taux d'actualisation
est de $3{,}50\,\%$. Calculer la valeur actuelle de cet échéancier.
\end{exercice}

\begin{corrige}
\begin{center}\small
\begin{tabular}{ccrr}
\toprule
Année & Flux & Facteur $v^{k}$ & Flux actualisé \\
\midrule
1 & $12\,000$ & $0{,}96618$ & $11\,594{,}20$ \\
2 & $15\,000$ & $0{,}93351$ & $14\,002{,}33$ \\
3 & $18\,000$ & $0{,}90194$ & $16\,234{,}86$ \\
4 & $20\,000$ & $0{,}87144$ & $17\,428{,}84$ \\
\midrule
\multicolumn{3}{r}{\textbf{Total}} & $\mathbf{59\,260{,}23}$ \\
\bottomrule
\end{tabular}
\end{center}
La valeur actuelle est d'environ $59\,260$~\euro, contre $65\,000$~\euro{} en
somme nominale. L'écart de $5\,740$~\euro{} est l'effet d'actualisation : c'est
lui, et lui seul, qui expliquera plus tard une part importante de la différence
entre une provision en normes françaises et une meilleure estimation.
\end{corrige}

\subsection{Généralisation : rentes certaines et taux de rendement}

\begin{generalisation}{Deux raccourcis et une lecture inverse}
\textbf{La rente certaine.} Lorsque tous les flux sont égaux à $C$ et versés en
fin de chacune des $n$ périodes, la somme admet une forme fermée :
\[
  V_0 \;=\; C \times a_{\overline{n}|r},
  \qquad
  a_{\overline{n}|r} \;=\; \frac{1-(1+r)^{-n}}{r}.
\]
Le coefficient $a_{\overline{n}|r}$\index{annuité} se lit comme le « nombre
d'années équivalent » que représente la rente en valeur d'aujourd'hui.

\medskip
\textbf{Le taux de rendement interne.} On peut inverser la question : quel taux
annule la VAN ? Ce taux $\rho$, appelé \emph{taux de rendement interne} ou
\emph{taux actuariel}\index{taux actuariel}, est solution de
\[
  -\,C_0 + \sum_{k=1}^{n} \frac{C_k}{(1+\rho)^{t_k}} \;=\; 0 .
\]
Il n'a pas d'expression analytique au-delà de deux périodes et s'obtient par
approximations successives.

\medskip
\textbf{L'absence d'opportunité d'arbitrage.} Deux échéanciers produisant
exactement les mêmes flux aux mêmes dates ont nécessairement le même prix. Sans
quoi il suffirait d'acheter le moins cher et de vendre le plus cher pour
encaisser un gain sans risque et sans mise de fonds. Ce principe, dit
d'\emph{absence d'opportunité d'arbitrage}\index{absence d'opportunité d'arbitrage},
est ce qui autorise à valoriser un engagement par ses flux plutôt que par son
prix de marché --- lequel n'existe pas, s'agissant d'un contrat d'assurance.
\end{generalisation}

\begin{exercice}[rente certaine et taux actuariel]
\textbf{(a)} Calculer la valeur actuelle d'une rente de $8\,000$~\euro{} par an
versée pendant $10$~ans, au taux de $2{,}80\,\%$.

\medskip
\textbf{(b)} Un investissement de $50\,000$~\euro{} rapporte $14\,000$~\euro{}
par an pendant $5$~ans. Déterminer son taux de rendement interne.
\end{exercice}

\begin{corrige}
\textbf{(a)} $(1{,}028)^{-10} = 0{,}75870$, donc
\[
  a_{\overline{10}|2{,}8\%} = \frac{1-0{,}75870}{0{,}028} = 8{,}6179
  \quad\Longrightarrow\quad
  V_0 = 8\,000 \times 8{,}6179 \approx 68\,943\ \text{\euro}.
\]

\medskip
\textbf{(b)} On cherche $\rho$ tel que $a_{\overline{5}|\rho} = 50\,000/14\,000 = 3{,}5714$.
Deux essais encadrent la solution :
\[
  a_{\overline{5}|12\%} = 3{,}6048
  \qquad
  a_{\overline{5}|13\%} = 3{,}5172 .
\]
Une interpolation linéaire entre ces deux points donne
$\rho \approx 12{,}0 + \frac{3{,}6048-3{,}5714}{3{,}6048-3{,}5172} \approx 12{,}4\,\%$.

L'interpolation linéaire n'est ici qu'une approximation : la fonction est
convexe, et le résultat exact se situe légèrement en deçà.
\end{corrige}

\begin{exercice}[d'un contrat décrit en mots à un échéancier]
Un contrat de capitalisation prévoit le versement d'une prime unique de
$100\,000$~\euro{} à la souscription, puis le service de $25\,000$~\euro{} en fin
de chacune des cinq années suivantes. Le taux de marché est de $4{,}00\,\%$.

Construire l'échéancier, calculer la VAN, et dire ce que l'on peut en conclure
sur le prix de ce contrat.
\end{exercice}

\begin{corrige}
Échéancier : $-100\,000$ en $t=0$, puis $+25\,000$ en $t=1,\dots,5$.
\[
  a_{\overline{5}|4\%} = \frac{1-(1{,}04)^{-5}}{0{,}04} = 4{,}45182
  \quad\Longrightarrow\quad
  \mathrm{VAN} = -100\,000 + 25\,000 \times 4{,}45182 = +11\,295{,}6\ \text{\euro}.
\]
La VAN est positive : au taux de marché retenu, le contrat vaut $111\,296$~\euro{}
alors qu'il se paie $100\,000$~\euro.

Par absence d'opportunité d'arbitrage, un tel écart ne peut pas se maintenir sur
un marché liquide : il signale soit que le taux de $4\,\%$ n'est pas le bon taux
d'actualisation pour cet engagement, soit que le contrat comporte un élément non
modélisé --- un risque de défaut, une option de rachat, une condition de survie.
Chacune de ces trois pistes fait l'objet d'une unité ultérieure.
\end{corrige}

\begin{pont}{Où cela sert dans Solvabilité~II}
Une \emph{cadence de règlement} en assurance non-vie n'est rien d'autre qu'un
échéancier : la charge ultime d'un exercice de survenance se règle sur plusieurs
années selon un profil observé, et la meilleure estimation des provisions pour
sinistres à payer est la valeur actuelle de cet échéancier. La formule employée
sera littéralement celle de cette unité.

Le passage d'une provision en normes françaises, non actualisée, à une meilleure
estimation actualisée produit un écart dont on a vu ci-dessus la mécanique :
c'est le premier des deux grands facteurs de divergence entre bilan comptable et
bilan prudentiel.
\end{pont}

\begin{refreg}
\textbf{Références.} Directive 2009/138/CE, article~77 : la meilleure estimation
correspond à la moyenne pondérée par leur probabilité des flux de trésorerie
futurs, compte tenu de la valeur temporelle de l'argent. \\
Règlement délégué (UE) 2015/35 : dispositions relatives à la projection des flux
de trésorerie et à la frontière des contrats. \emph{[Références d'articles à
préciser avant diffusion.]}
\end{refreg}

\subsection*{Notations introduites}
\begin{description}[leftmargin=2.2cm, style=nextline, font=\normalfont]
  \notation{C_k}{flux de la période $k$}
  \notation{\mathrm{VAN}}{valeur actuelle nette}
  \notation{a_{\overline{n}|r}}{valeur actuelle d'une rente certaine de 1 par an sur $n$ ans}
  \notation{\rho}{taux de rendement interne, ou taux actuariel}
\end{description}
